\documentclass[10pt,a4paper]{amsart}
\usepackage[margin=19mm]{geometry}
\usepackage{amsmath,amssymb,mathtools}
\usepackage[scr=rsfs,cal=euler]{mathalfa}
\usepackage{xfrac}
\usepackage[unicode,hidelinks,bookmarks=false]{hyperref}
\newcommand{\ie}{i.e.\ }
\newcommand{\cf}{cf.\ }
\newcommand{\resp}{respectively\ }
\newcommand{\Subsection}{\subsection}
\providecommand{\nobreakdash}{\nobreak\hbox{-}}
\setlength{\emergencystretch}{2em}
\multlinegap0pt
\usepackage{tikz}
\newtheorem{lemma}{Lemma}
\newtheorem{theorem}{Theorem}
\title{The finitude of the fibers of the complementary Bell numbers}
\author{John M. Campbell}
\date{Source: arXiv:2608.00575v1 (CC BY 4.0)}
\begin{document}
\maketitle

\noindent\emph{Source and licence.} The finitude of the fibers of the complementary Bell numbers. Version arXiv:2608.00575v1 (CC BY 4.0), distributed by arXiv under the Creative Commons Attribution 4.0 International licence, http://creativecommons.org/licenses/by/4.0/, as recorded in the arXiv licence metadata for this exact version and shown on the abstract page https://arxiv.org/abs/2608.00575v1. Full licence notice: \texttt{licenses/arxiv-2608.00575-CC-BY-4.0.txt}.\par\smallskip

\noindent\textit{Editorial scope.} Counted unit: the article's theorem maintheorem (source0.tex lines 687--690). The article's complete original proof of it is delivered with it (source0.tex lines 692--732, one \texttt{proof} environment of 41 lines). No mathematics is written, repaired or re-derived: the statement and the proof are the article's own lines, in the article's own notation, and no retained line is substituted or re-flowed.  Retained uncounted as the source-local context the proof needs: lines 506--509; lines 653--656.  Nothing was omitted from any retained span, so the package carries no omission record.  No retained line contains a citation, so the wrapper carries no bibliography and no bibliography entry.  No retained line contains a figure environment, an image-inclusion command or drawing code, so no image is delivered and the package declares no asset dependency.  Editorial formatting only, in addition: an AMS \texttt{amsart} preamble that restates the article's own declarations and macros used by the retained lines, namely the environments \texttt{lemma}, \texttt{theorem} on separate counters, together with the standard packages those lines need.  The retained environments print in this document as Lemma 1, Theorem 1 in order of appearance, whereas the article prints the same items with the numbering of its own full document.  The complete native source of the article as distributed in the arXiv e-print for this exact version is bundled unchanged as \texttt{sources/206547/source0.tex}.
\begin{lemma}\label{lemmaFf}
 For each $a \in \{ 0, 1, \ldots, 11 \}$   there is a power series $F_{a}(z)$ in $\mathbb{Q}_{2}\langle z \rangle$   such that:   For all $m \in  
  \mathbb{N}_{0}$, the equality   $F_{a}(z) |_{z = m} = f(a + 12 m)$ holds. 
\end{lemma}

\begin{lemma}\label{lemmainterpolate}
 For every $a \in \{ 0, 1, \ldots, 11 \}$, the interpolating function $F_{a}$ 
 is nonconstant. 
\end{lemma}

\begin{theorem}\label{maintheorem}
 For each $c \in \mathbb{Z}$, the fiber
$ \left\{ n \in \mathbb{N}_{0} : f(n) = c \right\} $ is finite. 
\end{theorem}

\begin{proof}
 Let $c$ be a fixed integer. For each value $ a$ in $ \{ 0, 1, \ldots, 11 \}$, 
 we define the $2$-adic analytic function 
\begin{equation}\label{defineGac} 
 G_{a, c}(z) := F_{a}(z) - c. 
\end{equation}
 By Lemma \ref{lemmainterpolate}, 
 the interpolating function $F_{a}$ is nonconstant. 
 This allows us to conclude that 
 the element of the Tate algebra given on either side of the equality in \eqref{defineGac} 
 is not the zero power series in $\mathbb{Q}_{2}\langle z \rangle$. 
 Strassmann's theorem then allows us to conclude that the zero set
\begin{equation}\label{defineZac}
 Z_{a, c} := \{ z \in \mathbb{Z}_{2} : G_{a, c}(z) = 0 \} 
\end{equation}
 is finite. 
 Now, let $m \in \mathbb{N}_{0}$. 
 Lemma \ref{lemmaFf} then gives us the biconditional equivalences such that 
\begin{equation}\label{biconditional}
 f(a+12m) = c 
 \ \Longleftrightarrow 
 \ F_{a}(m) - c = 0 
 \ \Longleftrightarrow \ m \in Z_{a, c} \subseteq \mathbb{Z}_{2}
\end{equation}
 So, from the inclusion of $\mathbb{N}_{0}$ in $\mathbb{Z}_{2}$, 
 along with the finitude the family defined in \eqref{defineZac}, together with the equivalences in 
 \eqref{biconditional}, we find that the set 
\begin{equation}\label{firstfinite}
 \{ m \in \mathbb{N}_{0} : f(a + 12 m) = c \} 
\end{equation}
 is finite: 
 By then using the decomposition of $\mathbb{N}_{0}$ into residue classes modulo $12$, 
 we find that 
\begin{equation}\label{secondfinite} 
 \{ n \in \mathbb{N}_{0} : f(n) = c \} 
 = \bigcup_{a=0}^{11} \{ a + 12 m : m \in \mathbb{N}_{0}, f(a + 12 m) = c \}.
\end{equation}
 From the finitude of \eqref{firstfinite}, 
 the right-hand side of \eqref{secondfinite} 
 is a finite union of finite sets, and is, consequently, finite. 
\end{proof}

\end{document}
